% !PRINTED 220
\weberpage{220}
%<<CONT
Octa\"edersubstitutionen absolut unge\"andert, w\"ahrend $K$ wieder
die gleichen Vorzeichen\"anderungen wie $f$ erleidet.

\begin{center}
\S.\,58.

\gesperrt{Die Ikosa\"edergruppe.}
\end{center}

Da wir bei der Ikosa\"edergruppe nur ein System conjugirter
f\"unfz\"ahliger Pole haben, so k\"onnen wir (\S.\,53, 1.) in dieser Gruppe
die beiden Substitutionen

\weq{1}{\Theta(x) = \varepsilon x, \quad \psi(x) = \frac{-1}{x}}

annehmen, worin $\varepsilon$ eine primitive f\"unfte Einheitswurzel bedeutet.
Wir nehmen hier, was freisteht, die Substitutionen $\psi$ in der
Form $-1:x$ an, weil dadurch die Formeln einfacher werden.
Die zu dem Systeme der f\"unfz\"ahligen Pole geh\"orige Grundform
12\textsuperscript{ten} Grades muss, da sie die Substitutionen (1) gestattet, von
der Form sein:

\weq{2}{f(x_1,x_2) = x_1x_2(x_1^{10} + m\,x_1^5x_2^5 - x_2^{10}),}

und es handelt sich noch um die Bestimmung des constanten
Factors $m$. Die beiden anderen Grundformen sind vom 20\textsuperscript{sten}
und 30\textsuperscript{sten} Grade. Wir k\"onnen nun $m$ nach dem Satze \S.\,54, 4.
bestimmen, wenn wir eine Covariante von $f(x)$ bilden k\"onnen,
deren Grad niedriger als 60 ist, und die sich nicht als ein Product
 aus Potenzen von drei Functionen 12\textsuperscript{ten}, 20\textsuperscript{sten}, 30\textsuperscript{sten}
Grades darstellen l\"asst, die nach dem erw\"ahnten Satze identisch
Null sein muss.

Nun l\"asst sich leicht eine Covariante 16\textsuperscript{ten} Grades von der
Form $f$ bilden, wenn wir nach Bd.\ I, \S.\,60 die vierte Polare der
Form $f(x_1,x_2)$ nehmen:

\[
12\,.\,11\,.\,10\,.\,9\;P_4(x,\xi) = u_0\xi_1^4 + u_1\xi_1^3\xi_2 + u_2\xi_1^2\xi_2^2 + u_3\xi_1\xi_2^3 + u_4\xi_2^4,
\]

worin

\[
u_0 = \frac{\partial^4 f}{\partial x_1^4}, \quad u_1 = 4\,\frac{\partial^4 f}{\partial x_1^3\,\partial x_2}, \quad u_2 = 6\,\frac{\partial^4 f}{\partial x_1^2\,\partial x_2^2},
\]
\[
u_3 = 4\,\frac{\partial^4 f}{\partial x_1\,\partial x_2^3}, \quad u_4 = \frac{\partial^4 f}{\partial x_2^4}.
\]

Daraus erhalten wir eine Covariante 16\textsuperscript{ten} Grades von $f$ als
%CONT>>
